\documentclass[11pt,letterpaper]{article}

\usepackage[margin=1in]{geometry}
\usepackage[T1]{fontenc}
\usepackage[utf8]{inputenc}
\usepackage{lmodern}
\usepackage{microtype}
\usepackage{parskip}
\usepackage{hyperref}
\usepackage{enumitem}
\usepackage{xcolor}

\hypersetup{
  colorlinks=true,
  linkcolor=black,
  urlcolor=blue!60!black,
  citecolor=black,
}

\setlist[itemize]{leftmargin=1.4em,itemsep=2pt,topsep=3pt,parsep=0pt}

\title{}
\author{}
\date{}

\begin{document}

\noindent
{\Large\bfseries Revision Cover Letter\par}
\vspace{0.8em}

\noindent\textbf{Date:} April 2026

\noindent\textbf{To:} Editorial Office, \emph{Array} (Elsevier)

\noindent\textbf{Re:} Major revision of ``CLOP-DiT: Structured-Metadata-Conditioned Single-Cell Latent
Generation via Contrastive Language--Omics Pretraining and Diffusion Transformers''

\noindent\textbf{Authors:} Zeyu Fu$^{1,*}$, Yiyao Liu$^{2}$, Junping Wang$^{3}$, Song Wang$^{4}$

\vspace{0.6em}
\noindent\textbf{Corresponding author:}

\noindent Zeyu Fu\\
State Key Laboratory of Trauma and Chemical Poisoning, Institute of Combined Injury,\\
Chongqing Engineering Research Center for Nanomedicine, College of Preventive Medicine,\\
Army Medical University, Chongqing 400038, China\\
Email: \href{mailto:fuzeyu09@gmail.com}{fuzeyu09@gmail.com}

\vspace{0.8em}
\hrule
\vspace{0.8em}

\noindent Dear Editors,

Thank you for the opportunity to revise our manuscript. The reviewers' comments helped us sharpen
both the scientific framing and the supporting evidence, and we are submitting a substantially
revised version together with a marked manuscript and a detailed point-by-point response.

\section*{Summary of Revision Work}

In the revised manuscript, we addressed the reviewers' concerns in four main ways.

\subsection*{1. Clarified the scope and claims}

\begin{itemize}
  \item We revised the abstract and discussion to state more clearly what CLOP-DiT is intended to
    demonstrate: controllable structured-metadata-conditioned generation, rather than a
    high-fidelity substitute for real single-cell data.
  \item We softened wording that was too strong or too vague, including claims about template
    optimality, ``understanding'' of biological semantics, and the interpretation of rare-cell
    failure modes.
  \item We made the distinction between stage-specific CLOP alignment metrics and end-to-end
    generation quality more explicit.
\end{itemize}

\subsection*{2. Added targeted analyses that directly answer reviewer concerns}

\begin{itemize}
  \item We added a latent-space KNN error taxonomy, showing that 49.2\% of mis-typings remain
    within the same biological family rather than crossing into unrelated lineages.
  \item We added organism-stratified evaluation and found no statistically significant
    human-only vs.\ mouse-only difference on the headline metrics.
  \item We replaced the vague phrase ``biologically difficult'' with a quantitative analysis
    showing that within-type heterogeneity is a stronger predictor of fidelity than training-set
    abundance.
  \item We expanded the variance analysis so that residual plots and variance-ratio summaries are
    interpreted explicitly rather than left implicit.
  \item We added a rare-cell failure analysis showing that the main problem is latent
    under-dispersion, not decoder collapse.
\end{itemize}

\subsection*{3. Added targeted experiments where the reviewers asked for stronger evidence}

\begin{itemize}
  \item We added Gaussian variance-recovery baselines to contextualize the near-zero
    cross-dataset $r_{\mathrm{var}}$ values.
  \item We added a formal ZCA whitening ablation and found that ZCA provides a modest but
    measurable benefit rather than a dramatic dependence.
  \item We added an augmentation-strategy sweep with a forced-scarcity control, which supports a
    narrower conclusion: simple mixing strategies do not rescue the current generator under the
    natural regime, but augmentation becomes useful under genuine scarcity.
  \item We added a bridge experiment from CLOP-stage ablations to end-to-end generation and found
    that Stage-1 rankings should not be treated as reliable proxies for final generation quality.
\end{itemize}

\subsection*{4. Extended the manuscript beyond the original in-distribution evaluation}

\begin{itemize}
  \item We added a matched-dimensionality encoder comparison showing that the severe loss of
    within-type variation is specific to the scGPT encoder, not a universal property of
    dimensionality reduction.
  \item We added a zero-shot strict-OOD evaluation on kidney, cerebellum, and fetal gonadal
    tissue. The model generalizes partially to structurally familiar programs, but not to
    structurally distinct neuronal programs, which makes the current limit clearer and more
    defensible.
  \item We regenerated Figures~1 and~2 to improve readability and adjusted crowded visual
    elements in the results figures.
\end{itemize}

\section*{Key Manuscript Changes}

\begin{itemize}
  \item \textbf{Abstract} (Comment~2.1): Added research-background framing and challenge
    statement for label-conditioned generation.
  \item \textbf{Results} (Section~3.4): Integrated KNN error taxonomy, organism-stratified
    evaluation, and abundance--fidelity correlation analysis into the Per-Type Generation
    Quality section.
  \item \textbf{Gene-level fidelity} (Section~3.5): Added variance and residual
    interpretation and embedded Gaussian variance-recovery baselines directly in the text.
  \item \textbf{Ablation section} (Section~3.10, Table~6): Added formal
    ZCA ablation table and summarised the CLOP$\to$full-pipeline bridge experiment with
    an explicit partial-reversal caveat.
  \item \textbf{Rare-cell augmentation} (Section~3.11): Expanded failure-mechanism
    discussion and included the augmentation strategy sweep with forced-scarcity control.
  \item \textbf{New Results Section~3.13} \emph{Strict-OOD Generalization (Zero-Shot)}:
    Per-tissue and per-cell-type zero-shot evaluation on kidney, cerebellum, and fetal
    gonadal tissue (Figure~9 and Table~8 in the main text).
  \item \textbf{Discussion}: New unified mechanistic narrative (latent-stage under-dispersion
    is the bottleneck), encoder bottleneck diagnosis, Gaussian benchmark framing
    clarification, and an updated strict-OOD paragraph reflecting the new zero-shot
    evaluation.
  \item \textbf{Figures~1--2} (Comment~2.4): Regenerated with improved readability---larger
    fonts, wider component spacing, increased background contrast, and consistent line weights.
    Figure~1 now resolves the Shared-Space vs.\ $z_0$ overlap that was visible in the earlier
    draft.
  \item \textbf{New supplementary figures} supporting specific reviewer concerns: Figure~13
    (KNN family-confusion heatmap, Comment~2.6), Figure~14 (organism-stratified 6-panel
    boxplot with Mann--Whitney $p$-values, Comment~3.2), Figure~15 (forced-scarcity rescue
    curves across five augmentation strategies, Comment~3.3), and Figure~9 (per-cell-type
    zero-shot strict-OOD bar chart, Comment~3.1). Figures~9a and~9b (variance matching and
    gene--gene correlation) were also refined to shorten over-long annotation leader lines.
  \item \textbf{Language}: Replaced vague terms (``biologically difficult,''
    ``understands biological semantics,'' ``optimal five-field template'') with precise
    operational definitions throughout.
\end{itemize}

\section*{Accompanying Documents}

\begin{enumerate}
  \item \textbf{Revised manuscript}
  \item \textbf{Marked manuscript}
  \item \textbf{Point-by-point response letter} --- detailed responses to all 15 reviewer
    comments with original reviewer text quoted, quantitative evidence, and manuscript
    cross-references
  \item \textbf{This revision cover letter}
\end{enumerate}

\vspace{0.8em}
\noindent We believe the revised submission addresses the reviewers' concerns directly and presents
the contribution more precisely and more transparently than the original version. We appreciate the
editors' and reviewers' continued consideration.

\vspace{0.8em}
\noindent Sincerely,

\vspace{0.6em}
\noindent\textbf{Zeyu Fu} (corresponding author)\\
\href{mailto:fuzeyu09@gmail.com}{fuzeyu09@gmail.com}

\vspace{0.3em}
\noindent on behalf of Zeyu Fu, Yiyao Liu, Junping Wang, and Song Wang

\end{document}
